\section{Typing}
Judgments have the form \(\Gamma \vdash e : \tau\), where \(\Gamma\) is a standard typing context.

\paragraph{Core rules.}
\begin{mathpar}
  \inferrule[\textsc{Var}]
    {x{:}\tau \in \Gamma}
    {\Gamma \vdash x : \tau}
  \and
  \inferrule[\textsc{Abs}]
    {\Gamma, x{:}\tau_1 \vdash e : \tau_2}
    {\Gamma \vdash \lambda x.\,e : \tau_1 \to \tau_2}
  \and
  \inferrule[\textsc{Rec}]
    {\Gamma, f{:}\tau_1 \to \tau_2, x{:}\tau_1 \vdash e : \tau_2}
    {\Gamma \vdash \mathsf{rec}\;f\;x.\,e : \tau_1 \to \tau_2}
  \and
  \inferrule[\textsc{App}]
    {\Gamma \vdash e_1 : \tau_1 \to \tau_2 \\ \Gamma \vdash e_2 : \tau_1}
    {\Gamma \vdash e_1\,e_2 : \tau_2}
\end{mathpar}

\paragraph{Products and sums.}
\begin{mathpar}
  \inferrule[\textsc{Pair}]
    {\Gamma \vdash e_1 : \tau_1 \\ \Gamma \vdash e_2 : \tau_2}
    {\Gamma \vdash \langle e_1 , e_2\rangle : \tau_1 \times \tau_2}
  \and
  \inferrule[\textsc{Fst}]
    {\Gamma \vdash e : \tau_1 \times \tau_2}
    {\Gamma \vdash \fstkw\,e : \tau_1}
  \and
  \inferrule[\textsc{Snd}]
    {\Gamma \vdash e : \tau_1 \times \tau_2}
    {\Gamma \vdash \sndkw\,e : \tau_2}
  \and
  \inferrule[\textsc{Inl}]
    {\Gamma \vdash e : \tau_1}
    {\Gamma \vdash \mathsf{inl}\,e : \tau_1 + \tau_2}
  \and
  \inferrule[\textsc{Inr}]
    {\Gamma \vdash e : \tau_2}
    {\Gamma \vdash \mathsf{inr}\,e : \tau_1 + \tau_2}
  \and
  \inferrule[\textsc{Match}]
    {\Gamma \vdash e : \tau_1 + \tau_2 \\ \Gamma, x{:}\tau_1 \vdash e_1 : \tau \\ \Gamma, y{:}\tau_2 \vdash e_2 : \tau}
    {\Gamma \vdash \matchkw\;e\;\withkw\;\mathsf{inl}\,x \Rightarrow e_1 \mid \mathsf{inr}\,y \Rightarrow e_2 : \tau}
\end{mathpar}

\paragraph{Lists.}
\begin{mathpar}
  \inferrule[\textsc{Nil}]{\;}{\Gamma \vdash [\,] : \List{\tau}}
  \and
  \inferrule[\textsc{Cons}]
    {\Gamma \vdash e_1 : \tau \\ \Gamma \vdash e_2 : \List{\tau}}
    {\Gamma \vdash e_1 :: e_2 : \List{\tau}}
  \and
  \inferrule[\textsc{Match-List}]
    {\Gamma \vdash e : \List{\tau} \\ \Gamma \vdash e_1 : \sigma \\ \Gamma, x{:}\tau, xs{:}\List{\tau} \vdash e_2 : \sigma}
    {\Gamma \vdash \matchkw\;e\;\withkw\;[\,] \Rightarrow e_1 \mid x::xs \Rightarrow e_2 : \sigma}
\end{mathpar}

\paragraph{Booleans and control.}
\begin{mathpar}
  \inferrule[\textsc{True}]{\;}{\Gamma \vdash \mathsf{true} : \Bool}
  \and
  \inferrule[\textsc{False}]{\;}{\Gamma \vdash \mathsf{false} : \Bool}
  \and
  \inferrule[\textsc{If}]
    {\Gamma \vdash e : \Bool \\ \Gamma \vdash e_1 : \tau \\ \Gamma \vdash e_2 : \tau}
    {\Gamma \vdash \ifkw\;e\;\thenkw\;e_1\;\elsekw\;e_2 : \tau}
  \and
  \inferrule[\textsc{Let}]
    {\Gamma \vdash e_1 : \tau_1 \\ \Gamma, x{:}\tau_1 \vdash e_2 : \tau_2}
    {\Gamma \vdash \letkw\;x=e_1\;\inkw\;e_2 : \tau_2}
\end{mathpar}

\paragraph{Floats and arithmetic.}
\begin{mathpar}
  \inferrule[\textsc{FloatLit}]{\;}{\Gamma \vdash c : \Float{m}}
  \and
  \inferrule[\textsc{Neg}]
    {\Gamma \vdash e : \Float{m}}
    {\Gamma \vdash -e : \Float{m}}
  \and
  \inferrule[\textsc{Add}]
    {\Gamma \vdash e_1 : \Float{m_1} \\ \Gamma \vdash e_2 : \Float{m_2} \\ m_1 \preceq m \\ m_2 \preceq m}
    {\Gamma \vdash e_1 + e_2 : \Float{m}}
  \and
  \inferrule[\textsc{Mul-G}]
    {\Gamma \vdash e_1 : \Float{m_1} \\ \Gamma \vdash e_2 : \Float{m_2} \\ m_1 \preceq \G \\ m_2 \preceq \G}
    {\Gamma \vdash e_1 \times e_2 : \Float{\G}}
  \and
  \inferrule[\textsc{Mul-ConstL}]
    {\Gamma \vdash c : \Float{m_1} \\ \Gamma \vdash e : \Float{m_2} \\ m_1 \preceq m \\ m_2 \preceq m}
    {\Gamma \vdash c \times e : \Float{m}}
  \and
  \inferrule[\textsc{Mul-ConstR}]
    {\Gamma \vdash e : \Float{m_1} \\ \Gamma \vdash c : \Float{m_2} \\ m_1 \preceq m \\ m_2 \preceq m}
    {\Gamma \vdash e \times c : \Float{m}}
  \and
  \inferrule[\textsc{Cmp}]
    {\Gamma \vdash e_1 : \Float{m_1} \\ \Gamma \vdash e_2 : \Float{m_2} \\ m_1 \preceq \G \\ m_2 \preceq \G}
    {\Gamma \vdash e_1 < e_2 : \Bool}
\end{mathpar}

\paragraph{Distributions.}
\begin{mathpar}
  \inferrule[\textsc{Uniform}]
    {\Gamma \vdash e_1 : \Float{m_1} \\ \Gamma \vdash e_2 : \Float{m_2} \\ m_1 \preceq m \\ m_2 \preceq m}
    {\Gamma \vdash \uniform(e_1,e_2) : \Float{m}}
  \and
  \inferrule[\textsc{Gauss}]
    {\Gamma \vdash e_1 : \Float{m_1} \\ \Gamma \vdash e_2 : \Float{m_2} \\ m_1 \preceq m \\ m_2 \preceq \G}
    {\Gamma \vdash \gaussian(e_1,e_2) : \Float{m}}
  \and
  \inferrule[\textsc{Poisson}]
    {\Gamma \vdash e_1 : \Float{m_1} \\ m_1 \preceq m}
    {\Gamma \vdash \poisson(e_1) : \Float{m}}
  \and
  \inferrule[\textsc{Exponential}]
    {\Gamma \vdash e_1 : \Float{m_1} \\ m_1 \preceq \G}
    {\Gamma \vdash \exponential(e_1) : \Float{m}}
  \and
  \inferrule[\textsc{Beta}]
    {\Gamma \vdash e_1 : \Float{m_1} \\ \Gamma \vdash e_2 : \Float{m_2} \\ m_1 \preceq \G \\ m_2 \preceq \G}
    {\Gamma \vdash \betafn(e_1,e_2) : \Float{m}}
  \and
  \inferrule[\textsc{Gamma}]
    {\Gamma \vdash e_1 : \Float{m_1} \\ \Gamma \vdash e_2 : \Float{m_2} \\ m_1 \preceq m \\ m_2 \preceq \G}
    {\Gamma \vdash \gammafn(e_1,e_2) : \Float{m}}
\end{mathpar}

Multiplication is general-mode only unless one side is a (typed) literal, in which case it follows the literal's mode; addition is parametric in the mode; comparison is only allowed at general mode. Uniform sampling preserves the incoming mode; Gaussian sampling returns the incoming mode but requires a general-typed variance.
\paragraph{Modes and subtyping.}
We relate modes with a preorder \(m_1 \preceq m_2\):
\begin{mathpar}
  \inferrule[\textsc{Mode-Refl}]{\;}{m \preceq m}
  \and
  \inferrule[\textsc{Mode-GE}]{\;}{\G \preceq \E}
\end{mathpar}
Float subtyping follows modes: general floats can be used where expectation-mode floats are required, and otherwise floats are covariant in their mode while invariant in everything else.
\begin{mathpar}
  \inferrule[\textsc{Sub-Float-Refl}]{\;}{\Float{m} <: \Float{m}}
  \and
  \inferrule[\textsc{Sub-Float-GE}]{\;}{\Float{\G} <: \Float{\E}}
  \and
  \inferrule[\textsc{Sub-Float}]
    {m_1 \preceq m_2}
    {\Float{m_1} <: \Float{m_2}}
  \and
  \inferrule[\textsc{Subsumption}]
    {\Gamma \vdash e : \tau_1 \\ \tau_1 <: \tau_2}
    {\Gamma \vdash e : \tau_2}
\end{mathpar}

\paragraph{Structural subtyping (standard).}
Standard covariant/contravariant rules apply to composite types:
\begin{mathpar}
  \inferrule[\textsc{Sub-Refl}]{\;}{\tau <: \tau}
  \and
  \inferrule[\textsc{Sub-Prod}]
    {\tau_1 <: \tau_1' \\ \tau_2 <: \tau_2'}
    {\tau_1 \times \tau_2 <: \tau_1' \times \tau_2'}
  \and
  \inferrule[\textsc{Sub-Sum}]
    {\tau_1 <: \tau_1' \\ \tau_2 <: \tau_2'}
    {\tau_1 + \tau_2 <: \tau_1' + \tau_2'}
  \and
  \inferrule[\textsc{Sub-List}]
    {\tau <: \tau'}
    {\List{\tau} <: \List{\tau'}}
  \and
  \inferrule[\textsc{Sub-Arr}]
    {\tau_1' <: \tau_1 \\ \tau_2 <: \tau_2'}
    {\tau_1 \to \tau_2 <: \tau_1' \to \tau_2'}
\end{mathpar}

